\documentclass[10pt,a4paper]{amsart}
\usepackage[margin=19mm]{geometry}
\usepackage{amsmath,amssymb,mathtools}
\usepackage[scr=rsfs,cal=euler]{mathalfa}
\usepackage{xfrac}
\usepackage[unicode,hidelinks,bookmarks=false]{hyperref}
\newcommand{\ie}{i.e.\ }
\newcommand{\cf}{cf.\ }
\newcommand{\resp}{respectively\ }
\newcommand{\Subsection}{\subsection}
\providecommand{\nobreakdash}{\nobreak\hbox{-}}
\setlength{\emergencystretch}{2em}
\multlinegap0pt
\usepackage{bm}
\usepackage{bbm}
\usepackage{dsfont}
\usepackage{xspace}
\usepackage{cancel}
\usepackage{mathtools}
\usepackage{amsthm}
\makeatletter
\@ifundefined{theorem}{\newtheorem{theorem}{Theorem}}{}
\@ifundefined{lemma}{\newtheorem{lemma}[theorem]{Lemma}}{}
\makeatother
\newcommand\la{{\lambda}}
\newcommand{\N}{\mathbb{N}}
\newcommand{\Z}{\mathbb{Z}}
\title[Bochner-Riesz means on the Heisenberg group]{Bochner-Riesz means on the Heisenberg group}
\author{Detlef M\"uller, Lars Niedorf, Andreas Seeger, Betsy Stovall}
\date{Source: arXiv:2607.02189v1 (CC BY 4.0)}
\begin{document}
\maketitle

\noindent\emph{Source and licence.} Bochner-Riesz means on the Heisenberg group. Version arXiv:2607.02189v1 (CC BY 4.0), distributed under the Creative Commons Attribution 4.0 International licence, as recorded in the arXiv licence metadata for this exact version and shown on the abstract page https://arxiv.org/abs/2607.02189v1. Full rights notice: licenses/arxiv-2607.02189-CC-BY-4.0.txt.\par\smallskip

\noindent\textit{Editorial scope.} Counted unit: the Lemma whose source label is \texttt{lem:decay-away-2} in the article (source0.tex lines 4233--4252, source label \texttt{lem:decay-away-2}), delivered together with the article's own complete original proof of it (source0.tex lines 4254--4340, one \texttt{proof} environment of 87 lines). No mathematics is written, repaired or re-derived: statement and proof are the article's own text line for line. Required context retained uncounted: VPhi (lines 4208--4210); WPhi (lines 4224--4226). Every definition, display and label of those lines is retained unchanged. Omitted from this package and not redistributed: 1--4207, 4211--4223, 4227--4232, 4253--4253, 4341--4850. Nothing inside a retained excerpt is omitted or altered: every retained body is delivered exactly as the article's lines are written, so no omission receipt of any kind accompanies this package. Citations: the retained excerpts carry no citation command at all, so no bibliography entry of the article is transcribed and no reference is redistributed. Reference adaptation: none was needed and none was made; every reference inside the delivered unit resolves inside it, so no unresolved reference or citation is produced. Printed slips kept literal: no printed word, symbol, hypothesis or number of the counted statement or of its proof was altered. Layout and class: the article's own class is \texttt{amsart} with packages including amssymb, amsfonts, amsthm, amsmath, bbm, bm, xcolor, graphicx, eucal, enumitem, hyperref, subcaption; the wrapper is the lane's pinned \texttt{amsart} editorial wrapper at 10pt on A4, which restates the theorem-like environments the retained text needs and adds the article's own macro definitions that the retained text needs (\texttt{\textbackslash N}, \texttt{\textbackslash Z}, \texttt{\textbackslash la}), with extra wrapper packages (bm, bbm, dsfont, xspace, cancel, mathtools). The delivered numbering is the unit's own. Assets: no retained span contains a figure environment, a graphics-inclusion command, a PDF-page inclusion, a table or a TikZ picture; no image file is copied into this package, the package declares no asset dependency and no image file is redistributed with it. Licence: version arXiv:2607.02189v1 of this article is distributed under the Creative Commons Attribution 4.0 International licence, as recorded in the arXiv licence metadata for this exact version and shown on the abstract page https://arxiv.org/abs/2607.02189v1. A full rights notice is delivered at licenses/arxiv-2607.02189-CC-BY-4.0.txt. Characters: every retained excerpt is pure ASCII, so the delivered engine, XeLaTeX with its default Latin Modern text font, reports no missing character.
\begin{equation}\label{VPhi}
V( e^{i \lambda \tilde \Phi}) =i\la \Delta x^2 (s,\tau)e^{i \lambda \tilde \Phi}.
\end{equation}

\begin{equation}\label{WPhi}
W (e^{i \lambda  \tilde\Phi} )=4 i\la \Delta u(s,\tau)e^{i \lambda \tilde \Phi}.
\end{equation}

\begin{lemma}\label{lem:decay-away-2}
For every $N,M\in\N$, the kernel $K_{\lambda,s}^{k,\ell}$ can be written as a finite series
\[
K_{\lambda,s}^{k,\ell} = \sum_{\alpha\in \N^2 \times \Z\times \N^6} K_{\lambda,s,N,M,\alpha}^{k,\ell}
\]
of summands of the form
\begin{multline*}
K_{\lambda,s,N,M,\alpha}^{k,\ell}(x, u)
 = 
\lambda^{n + \frac{3}{2}-N-M} \int_0^\infty \int_0^\infty  e^{i \lambda \tilde \Phi} \,\partial_\sigma^{\alpha_1} (a_{\lambda}(s \sigma))\, \partial_\tau^{\alpha_2}(\eta_\ell(\tau-k \pi))\\
\times \left(\frac{\tau}{\sin \tau}\right)^{n} \frac{\sigma^{\alpha_3} \left(\sin \tau\right)^{\alpha_5} P_{\alpha_6}(\cos\tau,\sin\tau)}{\tau^{\alpha_ 4 } \left( \Delta x^2 (s,\tau) \right)^{\alpha_7}\left( \Delta u(s,\tau) \right)^{\alpha_8} } \, s^{\alpha_9} \, d\sigma \, d \tau,
\end{multline*}
where $P_{\alpha_6}\in \Z[X,Y]$ is a polynomial with integer coefficients and $\alpha = (\alpha_1,\dots,\alpha_9)\in \N^2 \times \Z\times \N^6$ satisfies the relations
\begin{itemize}
 \item $\alpha_ 4 -2\alpha_7-2\alpha_8\ge -N-M$,
 \item $\alpha_5-\alpha_2-2\alpha_7-\alpha_8\ge -N-M$,
 \item $\alpha_7\ge N$ and $\alpha_8\ge M$.
\end{itemize}
Moreover, if $N = 0$, then $\alpha_7 = 0$, and if $M = 0$, then $\alpha_8 = 0$.
\end{lemma}

\begin{proof}
We do induction over $|N| + |M|$. For $|N| + |M| = 0$ we have $K_{\lambda,s}^{k,\ell} = K_{\lambda,\alpha}^{k,\ell}$ if $P_{\alpha_6} = 1$ and $\alpha_j = 0$ for all $j$.

 Now suppose that the statement holds for the tuple $(N,M)$. We will show that the statement then also holds for $(N + 1,M)$ and $(N,M + 1)$. We use the notation
\begin{multline*}
K_{\lambda,s,N,M,\alpha}^{k,\ell}[E](x, u)
 = \lambda^{n + \frac{3}{2}-N-M} \int_0^\infty \int_0^\infty E(x,u,s,\tau,\sigma) \, \partial_\sigma^{\alpha_1} (a_{\lambda}(s \sigma)) \\
 \times \partial_\tau^{\alpha_2}(\eta_\ell(\tau-k \pi))
 \left(\frac{\tau}{\sin \tau}\right)^{n} \frac{\sigma^{\alpha_3} \left(\sin \tau\right)^{\alpha_5} P_{\alpha_6}(\cos\tau,\sin\tau)}{ \tau^{\alpha_ 4 } \left( \Delta x^2 (s,\tau) \right)^{\alpha_7}\left( \Delta u(s,\tau) \right)^{\alpha_8} } \, s^{\alpha_9} \, d\sigma \, d \tau.
\end{multline*}
With this notation, $K_{\lambda,s,N,M,\alpha} = K_{\lambda,s,N,M,\alpha}^{k,\ell}[e^{i\lambda \tilde \Phi}]$.
 By \eqref{VPhi}, this yields
\begin{align*}
& K_{\lambda,s,N,M,\alpha}^{k,\ell}[i\,e^{i\lambda \tilde \Phi}]
 = K_{\lambda,s,N + 1,M,\alpha + e_7}^{k,\ell}[V(e^{i\lambda \tilde \Phi})] \\ 
& = -K_{\lambda,s,N + 1,M,\alpha + 2e_ 4 + 2e_5 + e_7}^{k,\ell}[\partial_\sigma(e^{i\lambda \tilde \Phi})] + K_{\lambda,s,N + 1,M,\alpha-e_3 + e_ 4 + 2e_5 + e_7}^{k,\ell}[\partial_\tau(e^{i\lambda \tilde \Phi})].
\end{align*}
When integrating by parts  with respect to $\sigma$ in the first integral and employing the Leibniz  rule, we see that $\partial_\sigma$ either hits $\partial_\sigma^{\alpha_1} (a_{\lambda}(s \sigma))$, which increases $\alpha_1$, or $\sigma^{\alpha_3}$, which decreases $\alpha_3$. By induction hypothesis,
\begin{align*}
(\alpha_ 4 + 2)-2(\alpha_7 + 1)-2\alpha_8 & > -(N + 1)-M, \\
(\alpha_5 + 2)-\alpha_2-2(\alpha_7 + 1)-\alpha_8 & > -(N + 1)-M.
\end{align*}
For the second integral, when  integrating  by parts is $\tau$ and  employing the Leibniz   rule again, we make the following observations:
\begin{itemize}
\item[(i)] When $\partial_\tau$ hits $\partial_\tau^{\alpha_2}(\eta_\ell(\tau-k \pi))$, then $\alpha_2$ is additionally increased by 1. However, the induction hypothesis still implies
\[
(\alpha_5 + 2)-(\alpha_2 + 1)-2(\alpha_7 + 1)-2\alpha_8 \ge -(N + 1)-M.
\]
\item[(ii)] When $\partial_\tau$ hits $\tau^{n}$ or $1/\tau^{\alpha_ 4 }$, then $\alpha_ 4 $ is additionally increased by 1. This additional gain is not relevant for our purposes.

\item[(iii)] When $\partial_\tau$ hits $1/(\sin\tau)^{n}$ or $(\sin \tau)^{\alpha_5}$, then $\alpha_5$ changes additionally by  $-1$, but $\alpha_2$ is unchanged, and 
\[
(\alpha_5 + 1)-\alpha_2-2(\alpha_7 + 1)-\alpha_8  \ge -(N + 1)-M.
\]
\item[(iv)] When $\partial_\tau$ hits $P_{\alpha_6}(\cos\tau,\sin\tau)$, we just pass to another polynomial.

\item[(v)] When $\partial_\tau$ hits $(\Delta x^2 (s,\tau))^{-(\alpha_7+1)}$, we get
\[
\partial_\tau (\Delta x^2 (s,\tau))^{-(\alpha_7+1)}
 =( \alpha_7+1) (\Delta x^2 (s,\tau))^{-(\alpha_7 + 1)} \partial_\tau\Big( s \, \frac{\sin^2\tau}{\tau^2} \Big).
\]
Thus $\alpha_7$ is increased in total by 2. If $\partial_\tau$ hits the enumerator  of the factor $\sin^2\tau/\tau^2$, $\alpha_ 4 $ and $\alpha_5$ are both increased in total by 3, and if $\partial_\tau$ hits the denominator   of the factor $\sin^2\tau/\tau^2$, then both $\alpha_ 4 $ and $\alpha_5$ are increased in total by 4.   Note that
\begin{align*}
(\alpha_ 4 + 3 )-2(\alpha_7 + 2)-2\alpha_8 & \ge -(N + 1)-M, \\
(\alpha_5 + 3)-\alpha_2-2(\alpha_7 + 2)-\alpha_8 & \ge -(N + 1)-M,
\end{align*}
and if we increase both $\alpha_4$ and $\alpha_5$ by 4, we even get strict inequalities.

\item[(vi)] When $\partial_\tau$ hits $(\Delta u(s,\tau))^{-\alpha_8}$, we get
\[
\partial_\tau \left(\Delta u(s,\tau)\right)^{-\alpha_8}
 = \alpha_8 \left(\Delta u(s,\tau)\right)^{-(\alpha_8 + 1)} \partial_\tau\Big( s \, \frac{\tau-\sin\tau\cos\tau}{4\tau^2} \Big).
\]
Similar as in (v), $\alpha_8$ is increased by 1, while the factor $\partial_\tau( (\tau-\sin\tau\cos\tau)/\tau^2 )$ leads to two summands where $\alpha_ 4 $ or $\alpha_5$ is additionally increased.  If $\partial_\tau$ hits $1/\tau,$ or the enumerator of $-\sin\tau\cos\tau/\tau^2,$ then $\alpha_4$ is increased in total by $3$. Note here that
\begin{align*}
(\alpha_ 4 + 3 )-2(\alpha_7 + 1)-2(\alpha_8 + 1) &\ge -(N + 1)-M,\\
(\alpha_5 + 2)-\alpha_2-2(\alpha_7 + 1)-(\alpha_8 + 1) &\ge -(N + 1)-M.
\end{align*}
If $\partial_\tau$ hits  the denominator of $-\sin\tau\cos\tau/\tau^2,$ then $\alpha_4$ is increased in total by $4,$ and we even get strict inequalities.
\end{itemize}

This verifies the statement for $(N + 1,M)$. We can argue similarly for $(N,M + 1)$ by using \eqref{WPhi}. 
This yields
\begin{multline*}
 K_{\lambda,s,N,M,\alpha}^{k,\ell}[4i\,e^{i\lambda \tilde \Phi}]
 = K_{\lambda,s,N,M + 1,\alpha + e_8}[W(e^{i\lambda \tilde \Phi})] \\
 = K_{\lambda,s,N,M + 1,\alpha + e_ 4 + e_8}[\partial_\sigma(e^{i\lambda \tilde \Phi})]
-K_{\lambda,s,N,M + 1,\beta + 2e_ 4 + e_5 + e_8}[\partial_\sigma(e^{i\lambda \tilde \Phi})] \\
+ K_{\lambda,s,N,M + 1,\gamma - e_3 + e_ 4 + e_5 + e_8}[\partial_\tau(e^{i\lambda \tilde \Phi})]
\end{multline*}
with some new polynomials $P_{\beta_6},P_{\gamma_6}\in \Z[X,Y]$ and $\alpha_j = \beta_j = \gamma_j$ for all $j\neq 6$. For the first two summands, observe that
\begin{align*}
(\alpha_ 4 + 1)-2\alpha_7-2(\alpha_8 + 1) &\ge -N-(M + 1),\\
(\alpha_5 + 1)-\alpha_2-2\alpha_7-(\alpha_8 + 1) & > -N-(M + 1).
\end{align*}
For the third summand, we can argue similarly to (i)--(iv) above. For case (v), observe that
\begin{align*}
(\alpha_ 4 + 3)-2(\alpha_7 + 1)-2(\alpha_8 + 1) &\ge -N-(M + 1), \\
(\alpha_5 + 2)-\alpha_2-2(\alpha_7 + 1)-(\alpha_8 + 1) &\ge -N-(M + 1).
\end{align*}
For case (vi), observe that
\begin{align*}
(\alpha_ 4 + 3)-2\alpha_7-2(\alpha_8 + 2) &\ge -N-(M + 1), \\
(\alpha_5 + 1)-\alpha_2-2\alpha_7-(\alpha_8 + 2) &\ge -N-(M + 1).
\end{align*}
Altogether, we have verified the statement also for $(N,M+1),$ which  finishes the induction.
\end{proof}

\end{document}
